% 通用中英双栏模板；XeLaTeX 编译。默认输入类型为仅 PDF。
% ============================================================
% 论文中英对照译稿：排版参数与双栏环境
% 后续调整建议：先改“可调参数区”，通常不需要改环境实现。
% 编译方式：XeLaTeX；正文位于本文件后面的 \begin{document} 之后。
% ============================================================

% 文档类仍以五号初始化标题层级；正文精确字号由下方 BodyFontSize 统一设置。
% pair 环境继承正常正文；不使用 \small。
\documentclass[UTF8,fontset=fandol,zihao=5]{ctexart}

% ---------- 宏包：按功能分组 ----------
\usepackage{geometry}                         % 纸张与页边距
\usepackage{amsmath,amssymb,mathtools,amsthm} % 数学公式与符号
\usepackage{lmodern}                          % 西文字体
\usepackage{paracol}                          % 中英并排、独立流动的双栏
\usepackage{graphicx,xfp}\usepackage{booktabs,array,longtable}         % 表格
\usepackage{xcolor,fancyhdr,enumitem}         % 颜色、页脚与列表
\usepackage{needspace}                        % 为标题及后续文字预留空间
\usepackage{hyperref}              % 内部公式引用与链接

% ============================================================
% 一、可调参数区
% 单位提示：mm 适合页边距，pt 适合小间距，em 随当前字体大小变化。
% “plus / minus”允许 TeX 在分页时小幅伸缩间距。
% ============================================================

% ---------- 1. 页面与双栏 ----------
\newlength{\PageLeftMargin}   \setlength{\PageLeftMargin}{25mm}
\newlength{\PageRightMargin}  \setlength{\PageRightMargin}{25mm}
\newlength{\PageTopMargin}    \setlength{\PageTopMargin}{24mm}
\newlength{\PageBottomMargin} \setlength{\PageBottomMargin}{24mm}
\newlength{\BilingualColumnGap} \setlength{\BilingualColumnGap}{36mm}
\newlength{\BilingualRuleWidth} \setlength{\BilingualRuleWidth}{0.3pt}
\definecolor{ink}{HTML}{20313F}    % 正文标题色
\definecolor{accent}{HTML}{167D8D} % 次级标题与栏分隔线色
\newcommand{\BilingualRuleColor}{accent!35}
\newcommand{\DocumentLinkColor}{accent} % 链接主题色；可改为 ink 或其他已定义颜色

% ---------- 2. 正文行距、段距与换行 ----------
% 基础字号采用 bp（PDF 点，1 bp = 1/72 英寸）。英文、定理和独立公式均使用它。
\newcommand{\BodyFontBP}{11} % 英文基础字号的数值，单位 bp
\newcommand{\BodyFontSize}{\BodyFontBP bp}
% 数学正文随基础字号；上下标按原稿比例放大，避免文字与公式比例失衡。
\newcommand{\BodyMathScriptSize}{\fpeval{7*\BodyFontBP/10.5}pt}
\newcommand{\BodyMathSmallScriptSize}{\fpeval{5*\BodyFontBP/10.5}pt}
\newcommand{\BodyLineStretch}{1.3} % 沿用 ctex 默认行距倍率；如需更疏，可改为 1.4
% 输入类型只指用户实际提供的原文；MinerU 生成的公式源码不算 TeX 原文。
\newif\ifSourceHasTeX
\SourceHasTeXfalse % 有原 TeX 时改成 SourceHasTeXtrue
\newcommand{\EnglishIndentFactor}{1.5} % em；英文标题后的首段另设为零
\newcommand{\ChineseIndentFactor}{2} % 当前中文字宽；字号选定后才应用

\newlength{\BodyParagraphSkip}   \setlength{\BodyParagraphSkip}{4pt}
\newlength{\PairExtraGap}        \setlength{\PairExtraGap}{0pt}
% PairExtraGap 是每个 pair 结束后另外追加的间距；段距、公式间距另计。
% 默认 0pt，避免与公式上下间距叠加。建议只在确实需要时增加。
\newlength{\BodyEmergencyStretch} \setlength{\BodyEmergencyStretch}{1em}
% 以下两项只在含较高行内公式的行之间生效，不是通常的正文行距。
% 保留现有设置；若行内分式碰到上下行，可适当增加安全间隙。
\newcommand{\InlineMathSafetyGap}{\dimexpr\baselineskip-1em\relax}
\newlength{\InlineMathGapThreshold} \setlength{\InlineMathGapThreshold}{2.5pt}

% ---------- 3. 正文、定理内部的行间公式 ----------
% 下列默认值沿用当前正常字号的公式间距。
% 要调整普通段落或定理内部的公式，只修改这一组。
\newlength{\BodyDisplayAbove} \setlength{\BodyDisplayAbove}{10pt plus 2pt minus 5pt}
\newlength{\BodyDisplayBelow} \setlength{\BodyDisplayBelow}{10pt plus 2pt minus 5pt}
\newlength{\BodyDisplayShortAbove} \setlength{\BodyDisplayShortAbove}{0pt plus 3pt}
\newlength{\BodyDisplayShortBelow} \setlength{\BodyDisplayShortBelow}{6pt plus 3pt minus 3pt}
% “Short”用于公式前一行较短的情况，因此视觉间距可能略有差别。

% ---------- 4. 独立公式双栏块 equationpair ----------
\newlength{\EquationPairAboveGap} \setlength{\EquationPairAboveGap}{4pt}
\newlength{\EquationPairBelowGap} \setlength{\EquationPairBelowGap}{4pt}
% 这两个值是公式块的附加留白，不是文字基线到公式基线的总距离。
% equationpair 会自动补偿新栏开头的空基线；不要手动加负 \vspace。
% 不要在 equation、align 等行间公式环境前插入 \par 或空白段落。

% ---------- 5. 标题样式、间距及页末保护 ----------
\newcommand{\DocumentTitleStyle}{\LARGE}
\newcommand{\MainHeadingStyle}{\Large\bfseries\color{ink}}
\newcommand{\SubHeadingStyle}{\large\bfseries\color{accent}}
\newlength{\MainHeadingBefore} \setlength{\MainHeadingBefore}{8pt}
\newlength{\SubHeadingBefore}  \setlength{\SubHeadingBefore}{6pt}
% 以下整数表示标题前至少预留多少个正文行高，包含标题自身。
% 标题较长时可增加；它们不会把整个后续段落锁在同一页。
\newcommand{\MainHeadingKeepLines}{5}
\newcommand{\SubHeadingKeepLines}{4}
\newcommand{\StatementKeepLines}{3}

% ---------- 6. 定理、命题、引理和定义 ----------
\newlength{\StatementBeforeSkip} \setlength{\StatementBeforeSkip}{5pt}
\newlength{\StatementAfterSkip}  \setlength{\StatementAfterSkip}{5pt}
\newcommand{\StatementLabelStyle}{\bfseries}
\newcommand{\TheoremEnglishStyle}{\itshape}
\newcommand{\TheoremChineseStyle}{\itshape} % 当前 Fandol 配置下中文使用楷体
\newcommand{\DefinitionEnglishStyle}{\upshape}
\newcommand{\DefinitionChineseStyle}{\upshape}
% 中文定理若希望使用直立宋体，将 TheoremChineseStyle 改为 \upshape。

% ---------- 7. 页脚 ----------
\newcommand{\FooterTextStyle}{\normalsize}

% ============================================================
% 二、应用参数：将上面的配置连接到 LaTeX 排版长度
% ============================================================

% 纸张模式统一在这一行选择；当前为 A3 横版。
\geometry{a3paper,landscape,
  left=\PageLeftMargin, right=\PageRightMargin,
  top=\PageTopMargin, bottom=\PageBottomMargin}
\linespread{\BodyLineStretch}
\setlength{\columnsep}{\BilingualColumnGap}
\setlength{\columnseprule}{\BilingualRuleWidth}
\colseprulecolor{\BilingualRuleColor}
% 链接文字使用主题色；内部引用、文献、网址与文件链接保持一致。
\hypersetup{
  colorlinks=true,
  linkcolor=\DocumentLinkColor,
  citecolor=\DocumentLinkColor,
  urlcolor=\DocumentLinkColor,
  filecolor=\DocumentLinkColor
}
\setlength{\emergencystretch}{\BodyEmergencyStretch}

% 每个普通双栏块都显式采用同一套正文公式间距。
\newcommand{\ApplyBodyDisplaySpacing}{%
  \setlength{\abovedisplayskip}{\BodyDisplayAbove}%
  \setlength{\belowdisplayskip}{\BodyDisplayBelow}%
  \setlength{\abovedisplayshortskip}{\BodyDisplayShortAbove}%
  \setlength{\belowdisplayshortskip}{\BodyDisplayShortBelow}%
}
% 精确设置正常正文：基准行高为字号的 1.2 倍，再乘 BodyLineStretch。
% 因此只改 BodyFontSize 即可同时调整正常文字和行高。
\renewcommand{\normalsize}{%
  \fontsize{\BodyFontSize}{\dimexpr\BodyFontSize*6/5\relax}\selectfont
}
\DeclareMathSizes{\BodyFontSize}{\BodyFontSize}
  {\BodyMathScriptSize}{\BodyMathSmallScriptSize}
\AtBeginDocument{%
  \normalsize
  \setlength{\parindent}{0pt}%
  \setlength{\parskip}{\BodyParagraphSkip}%
  \setlength{\lineskip}{\InlineMathSafetyGap}%
  \setlength{\lineskiplimit}{\InlineMathGapThreshold}%
  \ApplyBodyDisplaySpacing
}

% ============================================================
% 三、双栏环境：英文左栏、中文右栏
% ============================================================

% 用法：\begin{pair} 英文 \switchcolumn 中文 \end{pair}
% 可选参数是块开始前预留的正文行数，如 \begin{pair}[3]。
% 参数默认为 0；长块仍可跨页。不能在 pair 内再嵌套 pair。
\newenvironment{pair}[1][0]{%
  \ifnum#1>0\relax
    \Needspace{#1\baselineskip}%
  \fi
  \begin{paracol}{2}\normalsize%
  \ApplyBodyDisplaySpacing
}{%
  \end{paracol}%
  \ifdim\PairExtraGap>0pt\relax
    \vspace{\PairExtraGap}%
  \fi
}

% 用法：两栏均只含行间公式时使用 equationpair。
% 普通段落或完整定理里的公式仍直接写在 pair 中，不另拆公式块。
% 这里的负间距是对新栏空基线的补偿；可调留白由上面的正值控制。
\newenvironment{equationpair}{%
  \begin{pair}%
  \setlength{\parskip}{0pt}%
  \ifdefined\setdisplayskipstretch
    \setdisplayskipstretch{1}% 若另行启用 setspace，避免重复放大公式间距
  \fi
  \setlength{\abovedisplayskip}{\dimexpr\EquationPairAboveGap-\baselineskip\relax}%
  \setlength{\abovedisplayshortskip}{\abovedisplayskip}%
  \setlength{\belowdisplayskip}{\EquationPairBelowGap}%
  \setlength{\belowdisplayshortskip}{\EquationPairBelowGap}%
}{%
  \end{pair}%
}

% ============================================================
% 四、双语标题：标题文字、字号和留白分别管理
% ============================================================

% 单栏标题组件：#1 为前间距，#2 为字体样式，#3 为标题文字。
\newcommand{\BPTAlignleft}{\raggedright}
\newcommand{\BPTAligncenter}{\centering}
\newcommand{\HeadingBlock}[4]{%
  \par\addvspace{#2}%
  \noindent{\csname BPTAlign#1\endcsname #3 #4\par}%
}
\newcommand{\BilingualHeading}[6]{%
  \begin{pair}[#1]
  \HeadingBlock{#4}{#2}{#3}{#5}\switchcolumn
  \HeadingBlock{#4}{#2}{#3}{#6}\end{pair}}
% 可选参数为 left / center；标题本身的编号使用原论文文字。
\NewDocumentCommand{\phead}{O{left} m m}{%
  \BilingualHeading{\MainHeadingKeepLines}{\MainHeadingBefore}{\MainHeadingStyle}{#1}{#2}{#3}}
\NewDocumentCommand{\psubhead}{O{left} m m}{%
  \BilingualHeading{\SubHeadingKeepLines}{\SubHeadingBefore}{\SubHeadingStyle}{#1}{#2}{#3}}
\NewDocumentCommand{\papertitle}{O{center} m m}{%
  \BilingualHeading{\MainHeadingKeepLines}{0pt}{\DocumentTitleStyle}{#1}{#2}{#3}}

% ============================================================
% 五、定理类环境：共享标签和留白逻辑，分别指定中英文字体
% ============================================================

% #1 为完整标签（例如 Theorem 7 / 定理 7），#2 为正文样式。
% 标签不自动编号，以便保留原论文编号及已有链接。
\newcommand{\StatementBegin}[2]{%
  \par\addvspace{\StatementBeforeSkip}%
  \noindent{\StatementLabelStyle #1.}\ % 标签后保留一个词间空格
  \begingroup#2%
}
\newcommand{\StatementEnd}{%
  \par\endgroup\addvspace{\StatementAfterSkip}%
}
\newenvironment{formaltheoremen}[1]{%
  \StatementBegin{#1}{\TheoremEnglishStyle}\ApplyEnglishIndent%
}{\StatementEnd}
\newenvironment{formaltheoremzh}[1]{%
  \StatementBegin{#1}{\TheoremChineseStyle}\ApplyChineseIndent%
}{\StatementEnd}
\newenvironment{formaldefinitionen}[1]{%
  \StatementBegin{#1}{\DefinitionEnglishStyle}\ApplyEnglishIndent%
}{\StatementEnd}
\newenvironment{formaldefinitionzh}[1]{%
  \StatementBegin{#1}{\DefinitionChineseStyle}\ApplyChineseIndent%
}{\StatementEnd}

% 一个完整定理的英文和中文陈述共用同一个 pair。
% 末尾公式必须位于相应 formal 环境及外层 pair 的结束标记之前。
% 定理块使用 pair[\StatementKeepLines]，只保护开头，不限制整块跨页。

% ============================================================
% 六、公式引用与页脚
% ============================================================

% \eqen{eq37} / \eqzh{eq37} 分别跳转到英文 / 中文栏同号公式。
% 标签约定：eq:eq37:en、eq:eq37:zh；正文已有标签不作改名。
% 带名称的编号引用：#1 为跳转目标，#2 为名称，#3 为编号。
% 例如 \NumberedLink{theorem:1:zh}{定理}{1}：名称黑色，编号主题色。
% 名称和编号仍组成同一个可点击链接；\textcolor 只改变名称颜色。
% 定理、命题、引理和定义自身的标题编号：黑色文字，无自链接。
% 跳转锚点仍由附近的 hypertarget 保留，正文引用继续使用 NumberedLink。
\newcommand{\StatementNumber}[2]{\textcolor{black}{#1~#2}}
\newcommand{\NumberedLink}[3]{%
  \hyperlink{#1}{\textcolor{black}{#2}~#3}%
}
% 用显式目标和原编号引用公式；只有确实存在于本文的对象才能作内部跳转。
\newcommand{\EqRef}[2]{(\hyperlink{#1}{#2})}
% 原脚注可用数字或符号；两栏均保留其原标记，内部目标由 paracol/hyperref 分离。
\newcommand{\SourceFootnote}[2]{\begingroup\renewcommand{\thefootnote}{#1}\footnote{#2}\endgroup}
% 原有 label 形式的公式引用也可使用下列组件。
\newcommand{\eqen}[1]{(\hyperref[eq:#1:en]{\ref*{eq:#1:en}})}
\newcommand{\eqzh}[1]{(\hyperref[eq:#1:zh]{\ref*{eq:#1:zh}})}
\pagestyle{fancy}
\fancyhf{}
% 不显示页眉及横线，仅在页脚中央保留页码。
\renewcommand{\headrulewidth}{0pt}
\renewcommand{\footrulewidth}{0pt}
\fancyfoot[C]{\FooterTextStyle\thepage}

% ============================================================
% 七、安全的普通段落配对与自动中文字号
% ============================================================

% ============================================================
% 中文正文：普通说明段落按高度自动选择中文字号
% ============================================================
% 普通中文段落自动调节；英文、完整定理和独立公式采用统一基础字号。
% 中文正文会从候选字号中自动选择最接近英文排版高度的一项。
% 达到字号上限后允许保留高度差，不强制放大到等高。
\usepackage{xparse}
\newcommand{\ChineseSizeOffsets}{0,0.25,0.5,0.75,1,1.25,1.5} % 单位 bp，按升序列出
\newcommand{\AdaptiveBaselineFactor}{1.2} % 基准行高/字号，再乘 BodyLineStretch
\newcommand{\AdaptiveMathTextSize}{\BodyFontSize} % 数学正文固定为本稿基础字号
\newcommand{\AdaptiveMathScriptSize}{\BodyMathScriptSize} % 继承统一上下标参数
\newcommand{\AdaptiveMathSmallScriptSize}{\BodyMathSmallScriptSize} % 二级上下标
% 基础字号变化时，数学字号按比例自动更新；需要改变上下标比例时再调整对应参数。
% 候选字号 = BodyFontBP + ChineseSizeOffsets；基本字号改动时数学字号同步变化。

\ExplSyntaxOn
\bool_new:N \l_bpt_first_bool
\bool_new:N \l_bpt_continued_bool
\bool_new:N \l_bpt_fixed_bool
\tl_new:N \l_bpt_id_tl
\int_new:N \l_bpt_keep_int
\keys_define:nn {bpt/textpair}
  {
    id .tl_set:N=\l_bpt_id_tl,
    mode .choice:,
    mode/auto .code:n={\bool_set_false:N \l_bpt_fixed_bool},
    mode/fixed .code:n={\bool_set_true:N \l_bpt_fixed_bool},
    first .bool_set:N=\l_bpt_first_bool,
    first .default:n=true,
    continued .bool_set:N=\l_bpt_continued_bool,
    continued .default:n=true,
    keep .int_set:N=\l_bpt_keep_int
  }
\cs_new_protected:Npn \ApplyEnglishIndent
  {
    \ifSourceHasTeX \setlength{\parindent}{\EnglishIndentFactor em}
    \else \setlength{\parindent}{0pt}\fi
    \bool_if:NT \l_bpt_first_bool {\setlength{\parindent}{0pt}}
    \bool_if:NT \l_bpt_continued_bool {\setlength{\parindent}{0pt}}
  }
\cs_new_protected:Npn \ApplyChineseIndent
  {
    \ifSourceHasTeX \setlength{\parindent}{\ChineseIndentFactor\ccwd}
    \else \setlength{\parindent}{0pt}\fi
    \bool_if:NT \l_bpt_continued_bool {\setlength{\parindent}{0pt}}
  }

\clist_new:N \l_bpt_adaptive_sizes_clist
\clist_new:N \l_bpt_offsets_clist
\clist_set:Nx \l_bpt_offsets_clist {\ChineseSizeOffsets}
\clist_map_inline:Nn \l_bpt_offsets_clist
  {\clist_put_right:Nx \l_bpt_adaptive_sizes_clist {\fp_eval:n{\BodyFontBP+(#1)}}}
% 为每个候选文字字号声明同一套数学字号。
\clist_map_inline:Nn \l_bpt_adaptive_sizes_clist
  {
    \DeclareMathSizes{#1bp}
      {\AdaptiveMathTextSize}
      {\AdaptiveMathScriptSize}
      {\AdaptiveMathSmallScriptSize}
  }

\seq_new:N \l_bpt_adaptive_columns_seq
\tl_new:N \l_bpt_adaptive_en_tl
\tl_new:N \l_bpt_adaptive_zh_tl
\tl_new:N \l_bpt_adaptive_chosen_size_tl
\box_new:N \l_bpt_adaptive_measure_box
\dim_new:N \l_bpt_adaptive_column_width_dim
\dim_new:N \l_bpt_adaptive_en_height_dim
\dim_new:N \l_bpt_adaptive_initial_zh_height_dim
\dim_new:N \l_bpt_adaptive_chosen_zh_height_dim
\dim_new:N \l_bpt_adaptive_candidate_height_dim
\fp_new:N \l_bpt_adaptive_best_error_fp
\fp_new:N \l_bpt_adaptive_candidate_error_fp
\int_new:N \g_bpt_adaptive_pair_int

% 设置中文候选字号；数学字号已在上面的 DeclareMathSizes 中固定。
\cs_new_protected:Npn \bpt_adaptive_font:n #1
  {
    \fontsize{#1bp}
      {\fp_eval:n { \AdaptiveBaselineFactor * (#1) } bp}\selectfont
    \ApplyBodyDisplaySpacing
    \setlength{\lineskip}{\InlineMathSafetyGap}
    \setlength{\lineskiplimit}{\InlineMathGapThreshold}\ApplyChineseIndent
  }

% 试排盒子仅用于测量，绝不作为最终正文输出，因此正文仍可正常跨页。
% #1 为字体设置代码，#2 为段落内容。
% 试排时取消链接操作，但保留显示文字，避免反复建立 PDF 链接。
% 自动试排仅适用于无标签、无脚注、无自动计数器的普通说明段落。
\cs_new_protected:Npn \bpt_adaptive_measure:nn #1#2
  {
    \vbox_set:Nn \l_bpt_adaptive_measure_box
      {
        \hsize=\l_bpt_adaptive_column_width_dim
        \linewidth=\l_bpt_adaptive_column_width_dim
        \textwidth=\l_bpt_adaptive_column_width_dim
        #1
        % 缩进由上面的语言字体选择器应用，试排与输出保持一致。
        \setlength{\parskip}{\BodyParagraphSkip}
        \cs_set_eq:NN \hyperlink \use_ii:nn
        \RenewDocumentCommand \hyperref { o m } { ##2 }
        #2\par
      }
  }

% 选择字号：以两栏自然排版高度之差的绝对值作为比较指标。
% 两个字号同样接近时，保留较小字号；仅一行的段落不调整。
\cs_new_protected:Npn \bpt_adaptive_choose_size:
  {
    \tl_set:Nx \l_bpt_adaptive_chosen_size_tl
      { \clist_item:Nn \l_bpt_adaptive_sizes_clist {1} }
    \bpt_adaptive_measure:nn {\normalsize\ApplyEnglishIndent}
      {\tl_use:N \l_bpt_adaptive_en_tl}
    \dim_set:Nn \l_bpt_adaptive_en_height_dim
      { \box_ht:N \l_bpt_adaptive_measure_box + \box_dp:N \l_bpt_adaptive_measure_box }
    \bpt_adaptive_measure:nn
      {\bpt_adaptive_font:n {\tl_use:N \l_bpt_adaptive_chosen_size_tl}}
      {\tl_use:N \l_bpt_adaptive_zh_tl}
    \dim_set:Nn \l_bpt_adaptive_initial_zh_height_dim
      { \box_ht:N \l_bpt_adaptive_measure_box + \box_dp:N \l_bpt_adaptive_measure_box }
    \dim_set_eq:NN \l_bpt_adaptive_chosen_zh_height_dim \l_bpt_adaptive_initial_zh_height_dim
    \fp_set:Nn \l_bpt_adaptive_best_error_fp
      { abs(\dim_to_fp:n {\l_bpt_adaptive_initial_zh_height_dim} -
            \dim_to_fp:n {\l_bpt_adaptive_en_height_dim}) }
    % 中文在基础字号下已经较高时不再放大；英文只有一行时也不调整。
    \dim_compare:nNnT {\l_bpt_adaptive_en_height_dim} > {\baselineskip}
      {
        \dim_compare:nNnT {\l_bpt_adaptive_initial_zh_height_dim} < {\l_bpt_adaptive_en_height_dim}
          {
            \clist_map_inline:Nn \l_bpt_adaptive_sizes_clist
              {
                \bpt_adaptive_measure:nn {\bpt_adaptive_font:n {##1}}
                  {\tl_use:N \l_bpt_adaptive_zh_tl}
                \dim_set:Nn \l_bpt_adaptive_candidate_height_dim
                  { \box_ht:N \l_bpt_adaptive_measure_box + \box_dp:N \l_bpt_adaptive_measure_box }
                \fp_set:Nn \l_bpt_adaptive_candidate_error_fp
                  { abs(\dim_to_fp:n {\l_bpt_adaptive_candidate_height_dim} -
                        \dim_to_fp:n {\l_bpt_adaptive_en_height_dim}) }
                \fp_compare:nNnT {\l_bpt_adaptive_candidate_error_fp} < {\l_bpt_adaptive_best_error_fp}
                  {
                    \fp_set_eq:NN \l_bpt_adaptive_best_error_fp \l_bpt_adaptive_candidate_error_fp
                    \tl_set:Nn \l_bpt_adaptive_chosen_size_tl {##1}
                    \dim_set_eq:NN \l_bpt_adaptive_chosen_zh_height_dim \l_bpt_adaptive_candidate_height_dim
                  }
              }
          }
      }
  }

% 正文仍使用英文 \switchcolumn 中文；可选参数改为明确的排版键。
% 此环境只支持一个 \switchcolumn，不在此环境中嵌套 pair。
% 测量使用当前等宽双栏的栏宽；修改 columnratio 时需同步修改此处。
% textpair[mode=auto|fixed, first, continued, keep=行数, id=块编号]
% 仅普通安全段落进入自动试排；间接包含副作用的宏必须显式 mode=fixed。
\NewDocumentEnvironment{textpair}{O{} +b}
  {
    \keys_set:nn {bpt/textpair}{id=,mode=auto,first=false,continued=false,keep=0,#1}
    \seq_set_split:Nnn \l_bpt_adaptive_columns_seq {\switchcolumn}{#2}
    \int_compare:nNnTF {\seq_count:N \l_bpt_adaptive_columns_seq}= {2}
      {
        \seq_pop_left:NN \l_bpt_adaptive_columns_seq \l_bpt_adaptive_en_tl
        \seq_pop_left:NN \l_bpt_adaptive_columns_seq \l_bpt_adaptive_zh_tl
        % 不展开任意宏，避免判断过程本身触发计数器或文件操作。
        \clist_map_inline:nn
          {footnote,footnotetext,footnotemark,SourceFootnote,label,hypertarget,
           phantomsection,stepcounter,refstepcounter,addtocounter,setcounter,
           caption,captionof,includegraphics,begin,write,index}
          {
            \exp_args:Nno \tl_if_in:nnT {#2}{\cs:w ##1\cs_end:}
              {\bool_set_true:N \l_bpt_fixed_bool}
          }
        \dim_set:Nn \l_bpt_adaptive_column_width_dim {(\textwidth-\columnsep)/2}
        \bool_if:NTF \l_bpt_fixed_bool
          {
            \tl_set:Nx \l_bpt_adaptive_chosen_size_tl {\BodyFontBP}
            \iow_term:x{CASE|\l_bpt_id_tl|fixed}
          }
          {
            \bpt_adaptive_choose_size:
            \int_gincr:N \g_bpt_adaptive_pair_int
            \iow_term:x{CASE|\l_bpt_id_tl|adaptive}
            \iow_term:x{ADAPTIVE|\int_use:N\g_bpt_adaptive_pair_int|
              \l_bpt_adaptive_chosen_size_tl|
              \dim_to_decimal:n{\l_bpt_adaptive_en_height_dim}|
              \dim_to_decimal:n{\l_bpt_adaptive_initial_zh_height_dim}|
              \dim_to_decimal:n{\l_bpt_adaptive_chosen_zh_height_dim}}
          }
        \begin{pair}[\int_use:N\l_bpt_keep_int]
        \normalsize\ApplyEnglishIndent\tl_use:N \l_bpt_adaptive_en_tl\par
        \switchcolumn
        \group_begin:
          \bpt_adaptive_font:n{\l_bpt_adaptive_chosen_size_tl}
          \tl_use:N \l_bpt_adaptive_zh_tl\par
        \group_end:
        \end{pair}
      }
      {\PackageError{bpt}{textpair~requires~exactly~one~switchcolumn}
        {Use~pair~for~structural~blocks.}}
  }{}
\ExplSyntaxOff


\begin{document}
\papertitle[center]{Bilingual Paper}{论文中英对照}
\phead[left]{1 Introduction}{1 引言}
\begin{textpair}[id=demo,first]
This is a compilable template demonstration. Replace this body with the paper's verified English blocks and their reviewed Chinese translations. Set SourceHasTeX from the actual input type, not from the extraction format.
\switchcolumn
这是可以直接编译的模板示例。生成译稿时，将示例正文替换为核验后的英文块及中文译文。输入类型依据用户提供的原文确定。
\end{textpair}
\end{document}
